\subsubsection{iCaRL}
\gls{icarl} \cite{ll_mof_2017} is an exemplar-based \gls{cil} method that jointly learns a feature representation and a classification model in an incremental fashion. It was originally proposed for \gls{cil}, where task identity is unavailable at inference time, and classification is performed using a nearest-mean-of-exemplars rule in the learned feature space. It has since been applied to \gls{til} settings in the literature, where the nearest-mean classifier is replaced by a standard linear head with task-indexed outputs. Therefore, \gls{icarl} is suitable for both \gls{cil} and \gls{til} settings.

\paragraph{Exemplar Management}
\gls{icarl} maintains a fixed total exemplar budget $K$ distributed across all observed classes.
For each class \(y\), it constructs an exemplar set \(\mathcal{P}_y=\{x_1,\ldots,x_m\}\) in feature space \(\phi(\cdot)\) to approximate the class mean:
\gls{icarl}'s memory is class-structured rather than task-structured: it stores exemplar inputs grouped by class label, typically as per-class sets \(\{\mathcal{P}_y\}_{y \in \mathcal{Y}}\), and uses these samples both for rehearsal and for prototype-based inference.
In the original \gls{cil} formulation, memory entries do not require explicit task identifiers; task identity is not assumed at test time, so prediction is performed over the active class set using class means derived from the stored exemplars.
Some practical codebases may additionally log or cache a task index per sample for bookkeeping, analysis, or \gls{til} evaluation, but this is an implementation convenience rather than a defining component of \gls{icarl}.
\begin{equation}
\mu_y=\frac{1}{|\mathcal{X}_y|}\sum_{x\in\mathcal{X}_y}\phi(x)
\end{equation}

Exemplars are selected greedily by herding. At step \(k\), the next exemplar is chosen from the remaining class samples as
\begin{equation}
x_k=\operatorname*{arg\,min}_{x\in \mathcal{X}_y\setminus\{x_1,\ldots,x_{k-1}\}}
\left\|
\mu_y-\frac{1}{k}\left(\phi(x)+\sum_{j=1}^{k-1}\phi(x_j)\right)
\right\|_2
\label{eqn:icarl_exemplar}
\end{equation}

When new classes are introduced, the fixed budget is redistributed over all observed classes. If \(|\mathcal{Y}|\) denotes the number of classes seen so far, the per-class quota is
\begin{equation}
    m=\left\lfloor \frac{K}{|\mathcal{Y}|}\right\rfloor
\end{equation}
and previously stored classes are truncated to at most \(m\) exemplars each. Thus, as \(|\mathcal{Y}|\) grows, older classes are progressively pruned.

\paragraph{\gls{cl} Approach}
\gls{icarl} mitigates \gls{cf} through two complementary mechanisms: experience replay using the stored exemplar sets, and knowledge distillation to preserve prior class representations. The combined training loss at task $t$ is:
%
\begin{equation}\label{eqn:icarl_loss}
    \mathcal{L} = \mathcal{L}_\text{cls}(x, y)
    + \lambda \sum_{\tau=0}^{t-1} |\mathcal{Y}_\tau| \cdot
      D_\text{KL}\!\left(
        \operatorname{softmax}\!\left(f_\theta(x^\tau)\right)
        \,\Big\|\,
        \operatorname{softmax}\!\left(f_{\theta^*_\tau}(x^\tau)\right)
      \right)
\end{equation}
%
where $\mathcal{L}_\text{cls}$ is the cross-entropy classification loss over the current task's classes $\mathcal{Y}_t$; $f_\theta(\cdot)$ denotes the logits of the current model; $f_{\theta^*_\tau}(\cdot)$ are the logits of the model as frozen at the end of task $\tau$, stored alongside the exemplars as distillation targets; $x^\tau$ is a sample drawn from the exemplar memory for task $\tau$; and $\lambda$ is a regularisation coefficient controlling the distillation penalty relative to the classification loss. The KL term penalises deviation of the current model's output distribution over old-class exemplars from the prior model's distribution, discouraging forgetting without requiring access to the original training data.

\paragraph{Inference}
In the \gls{cil} setting, \gls{icarl} classifies a test sample $x$ by assigning it to the class whose L2-normalised exemplar mean is nearest in feature space:
%
\begin{equation}\label{eqn:icarl_inference}
    \hat{y} = \operatorname*{arg\,min}_{y \in \mathcal{Y}}
    \left\|
      \frac{\phi(x)}{\|\phi(x)\|_2}
      - \frac{1}{\|\bar{\mu}_y\|_2}\bar{\mu}_y
    \right\|_2, \quad
    \bar{\mu}_y = \frac{1}{|\mathcal{P}_y|}\sum_{p \in \mathcal{P}_y} \phi(p)
\end{equation}
%
where both the query embedding and the per-class exemplar mean are normalised to the unit hypersphere before computing the distance, making the classifier equivalent to nearest-centroid under cosine similarity. The active class set $\mathcal{Y}$ is defined implicitly by the exemplar memory: only classes with stored exemplars are considered, making the classifier naturally suited to the \gls{cil} setting where the number of classes grows incrementally. In \gls{til} evaluations, the task identity is used to restrict the candidate set $Y$ to the classes of the current task.